% Part: history
% Chapter: biographies
% Section: rozsa-peter

\documentclass[../../../include/open-logic-section]{subfiles}

\begin{document}

\olfileid{his}{bio}{pet} 

\olsection{రోజా పీటర్}

\olphoto{peter-rozsa}{రోజా పీటర్}
 
రోజా పీటర్ 1905 ఫిబ్రవరి 17న హంగరీలోని బుడాపెస్ట్‌లో రోజా పొలిట్సర్‌గా జన్మించారు. పునరావృత్త ప్రమేయాలపై ఆమె చేసిన కృషికి ప్రసిద్ధి; పునరావృత్త సిద్ధాంతం అనే రంగం ఏర్పడటానికి ఆ కృషి అత్యవసరమైంది.

పీటర్ రాజకీయంగా కఠినమైన కాలంలో పెరిగారు---ఆమె కౌమార దశలో మొదటి ప్రపంచ యుద్ధం సాగింది---అయినా బుడాపెస్ట్‌లోని సంపన్న మారియా టెరెజియా బాలికల పాఠశాలలో చదివి, 1922లో ఉత్తీర్ణులయ్యారు. తరువాత బుడాపెస్ట్‌లోని పాజ్‌మానీ పీటర్ విశ్వవిద్యాలయంలో (తరువాత లోరాండ్ ఎట్వోష్ విశ్వవిద్యాలయంగా పేరు మారింది) చదివారు. తండ్రి పట్టుబట్టడంతో రసాయనశాస్త్రం చదవడం మొదలుపెట్టినా, తరువాత గణితానికి మారి 1927లో పట్టా పొందారు. ఉన్నత పాఠశాలలో గణితం బోధించడానికి అర్హతలున్నా, మహామాంద్యం ప్రపంచ ఆర్థిక వ్యవస్థను ప్రభావితం చేయడంతో అప్పటి ఆర్థిక పరిస్థితి దయనీయంగా ఉండేది. ఆ కాలంలో పీటర్ గణితంలో వ్యక్తిగత బోధకురాలిగా, ఇతర చిన్నచిన్న బోధనా పనులు చేశారు. తరువాత గణితంలో ఉన్నత అధ్యయనాల కోసం విశ్వవిద్యాలయానికి తిరిగి వెళ్లారు. మొదట సంఖ్యా సిద్ధాంతంలో పనిచేయాలనుకున్నారు; కానీ తాను కనుగొన్న ఫలితాలు ముందే నిరూపించబడ్డాయని తెలుసుకొని, గణితాన్నే పూర్తిగా వదిలేయబోయారు. గ్యోడెల్ అసంపూర్ణత సిద్ధాంతాలపై పనిచేయమని ప్రోత్సాహం పొందారు; ఆయన ఫలితాలలో కొన్నింటిని, అవి అవే ఫలితాలని తెలియకుండానే, వేరే మార్గాలలో నిరూపించారు. దీనితో ఆమెలో ఆత్మవిశ్వాసం తిరిగి పెరిగింది. డేవిడ్ హిల్బర్ట్ పునాదుల కార్యక్రమం స్ఫూర్తితో పునరావృత్త సిద్ధాంతంపై తన తొలి వ్యాసాలు రాశారు. 1935లో పీహెచ్‌డీ పొందారు; 1937లో \emph{Journal of
  Symbolic Logic} సంపాదకురాలయ్యారు.

పీటర్ తొలి వ్యాసాలను పునరావృత్త ప్రమేయ సిద్ధాంతం స్థాపనలో ముఖ్యమైన కృషిగా విస్తృతంగా గుర్తిస్తారు. \cite{Peter1935a}లో ఆమె వివిధ రకాల పునరావృత్తాల మధ్య సంబంధాన్ని పరిశీలించారు. \cite{Peter1935b}లో ఒక నిర్దిష్ట పునరావృత్తంగా నిర్వచించిన ప్రమేయం ఆదిమ పునరావృత్తం కాదని చూపారు. విల్హెల్మ్ ఆకెర్మాన్ సాధించిన పూర్వ ఫలితాన్ని ఇది సరళీకరించింది. పీటర్ సరళీకరించిన ప్రమేయాన్ని ఇప్పుడు తరచుగా ఆకెర్మాన్ ప్రమేయం అని, కొన్నిసార్లు మరింత సముచితంగా ఆకెర్మాన్--పీటర్ ప్రమేయం అని పిలుస్తారు. పునరావృత్త ప్రమేయ సిద్ధాంతంపై తొలి పుస్తకం కూడా ఆమె రాశారు \citep{Peter1951}.

ఆమె కృషి ముఖ్యమైనదీ ప్రభావవంతమైనదీ అయినప్పటికీ, 1945 వరకు పీటర్‌కు పూర్తి సమయ బోధనా పదవి లభించలేదు. రెండవ ప్రపంచ యుద్ధంలో హంగరీపై నాజీ ఆక్రమణ సమయంలో యూదు వ్యతిరేక చట్టాల వల్ల ఆమెకు బోధించడానికి అనుమతి లేకపోయింది. 1944లో ప్రభుత్వం బుడాపెస్ట్‌లో యూదుల ఘెట్టోను ఏర్పాటు చేసింది; దాన్ని మిగతా నగరం నుంచి వేరు చేసి సాయుధ రక్షకులను ఉంచారు. 1945లో విముక్తి పొందేవరకు పీటర్ ఆ ఘెట్టోలో నివసించవలసి వచ్చింది. తరువాత బుడాపెస్ట్ ఉపాధ్యాయ శిక్షణ కళాశాలలో, 1955 నుంచి ఎట్వోష్ లోరాండ్ విశ్వవిద్యాలయంలో బోధించారు. గణితంలో అకాడెమిక్ డాక్టర్ అయిన తొలి హంగేరియన్ మహిళా గణితశాస్త్రజ్ఞురాలు, హంగేరియన్ అకాడమీ ఆఫ్ సైన్సెస్‌కు ఎన్నికైన తొలి మహిళ కూడా ఆమె.

పీటర్ గణితాన్ని ఉత్సాహంగా బోధించేవారని పేరుంది. కేవలం ఉపన్యాసం ఇవ్వడం కంటే, విద్యార్థులతో కలిసి గణిత సమస్యల స్వభావం, సౌందర్యాన్ని అన్వేషించడానికి ఇష్టపడేవారు. అందువల్ల విద్యార్థులు ఆప్యాయంగా ఆమెను ``రోజా అత్తయ్య'' అని పిలిచేవారు. పీటర్ 1977లో~71 ఏళ్ల వయసులో మరణించారు.

\begin{reading}
మరిన్ని జీవిత చరిత్ర రచనలకు \citep{Oconnor2014}, \citep{Andrasfai1986} చూడండి. \citet{Tamassy1994} పీటర్‌తో సంక్షిప్త ముఖాముఖి నిర్వహించారు. గణితం గురించి సరదాగా చదవడానికి పీటర్ రాసిన \emph{Playing With Infinity} \citep{Peter2010} చూడండి.
\end{reading}

\end{document}
